\chapter{Problem Formulation}
\label{chap:problem_formulation}

This chapter formulates the responder-UAV planning problem used in the
Proactive Golden-Hour Bridge Model (PGBM). The formulation focuses on the
delivery layer after survivor reports have already been accepted by the
command center. Scout-UAV detection, human or AI-assisted verification, and
physical rescue by ground SAR teams are treated as external processes.

The responder planner decides which accepted survivor-assistance tasks should
be served in the current planning step, which UAV should serve each selected
task, which feasible three-dimensional route should be followed, and whether
the resulting missions satisfy payload and battery limits. The formulation
also includes a limited recourse rule for newly arriving tasks
when all responder UAVs are already committed.

\begin{figure}[H]
    \centering
    \includegraphics[width=0.86\textwidth]{pgbm_system_architecture.png}
    \caption{Overall PGBM system architecture showing scout-UAV detection,
    command-center task and UAV-state management, online responder assignment,
    energy-aware route planning, and mission-state feedback.}
    \label{fig:pgbm_system_architecture}
\end{figure}

\clearpage
\section{Main Notation}

Table~\ref{tab:main_notation} summarizes the main notation used throughout
the formulation. Key symbols are also briefly reintroduced near the
corresponding equations where needed for clarity.

\begin{longtable}{p{0.23\textwidth}p{0.67\textwidth}}
\caption{Main notation used in the PGBM formulation.}
\label{tab:main_notation}\\
\hline
\textbf{Symbol} & \textbf{Meaning}\\
\hline
\endfirsthead
\hline
\textbf{Symbol} & \textbf{Meaning}\\
\hline
\endhead
$\mathcal{I}(t)$ & Eligible survivor-assistance tasks at planning time $t$\\
$\mathcal{U}(t)$ & Available responder UAVs at planning time $t$\\
$\mathcal{R}$ & Relief-item types\\
$\mathcal{N}$ & Feasible UAV waypoints in the flyable 3D operating space\\
$\mathcal{E}$ & Feasible directed flight links between waypoints\\
$i,u,r$ & Task, UAV, and relief-item indices\\
$m,n$ & Waypoint indices\\
$x_{iu}$ & Binary task-assignment decision\\
$y_{mnu}$ & Binary flight-link selection decision\\
$s_i$ & Operational severity score of task $i$\\
$t_i$ & Detection time of task $i$\\
$\lambda$ & Common temporal decay rate of service value\\
$v_i(\tau)$ & Time-dependent service value of task $i$ at time $\tau$\\
$\delta_i$ & Delivery or service duration at task $i$\\
$b_u^{\mathrm{s}}, b_u^{\mathrm{r}}$ & Logical start and return base nodes of UAV $u$\\
$m_i$ & UAV-accessible drop-off waypoint for task $i$\\
$\mathbf{p}_m$ & 3D coordinate of waypoint $m$\\
$d_{mn}$ & 3D distance of flight link $(m,n)$\\
$t_{mnu}$ & Travel time of UAV $u$ on link $(m,n)$\\
$D_u,T_u$ & Total selected travel distance and travel time of UAV $u$\\
$\tau_{mu},C_i$ & Arrival time at waypoint $m$ and completion time of task $i$\\
$q_{ir}(t)$ & Remaining quantity of item type $r$ required by task $i$\\
$w_r$ & Weight of one item of type $r$\\
$K_u,W_u$ & Parcel-count and payload-weight capacities of UAV $u$\\
$M_i(t),M_u^0$ & Relief mass of task $i$ and initial payload of UAV $u$\\
$L_{mnu}$ & Payload mass carried by UAV $u$ on link $(m,n)$\\
$E_u(t),E_u^{\mathrm{res}}$ & Available and reserve battery energy of UAV $u$\\
$E_u^{\mathrm{req}}$ & Required mission energy of UAV $u$\\
$a_{mn}$ & Positive ascent on flight link $(m,n)$\\
$\alpha_u,\beta_u,\gamma_u$ & Energy coefficients for UAV $u$\\
$P_u^{\mathrm{srv}}$ & Service or hovering power of UAV $u$\\
$\mathcal{T}_u^{\mathrm{rem}}(t_r)$ & Uncompleted tasks of UAV $u$ at recourse time $t_r$\\
$J_{ui}^{\mathrm{old}},J_{ui}^{\mathrm{new}}$ & Remaining mission value before and after replacing task $i$\\
$\Delta J_{ui}$ & Objective gain from replacing task $i$\\
$\mathcal{S}(t_r),\mathcal{Q}(t_r)$ & Feasible replacement set and waiting task queue\\
\hline
\end{longtable}

\clearpage
\section{Decision Variables}

At planning time $t$, the command center considers a set of eligible tasks
$\mathcal{I}(t)$ and a set of available responder UAVs $\mathcal{U}(t)$. The
task-assignment decision is
\begin{equation}
x_{iu}
=
\begin{cases}
1, & \text{if task }i\text{ is assigned to UAV }u,\\
0, & \text{otherwise}.
\end{cases}
\end{equation}

The feasible flight environment is represented by a waypoint set
$\mathcal{N}$ and a directed flight-link set $\mathcal{E}$. A waypoint is a
candidate UAV position in flyable 3D space, and each waypoint $m$ has
coordinate $\mathbf{p}_m=(p_m^x,p_m^y,p_m^z)$. Buildings, rubble, unsafe
structures, and no-fly volumes are excluded when these waypoints and flight
links are constructed. A flight link $(m,n)$ is included in $\mathcal{E}$ only
if the straight segment between the two waypoints remains in flyable space.
The routing decision is
\begin{equation}
y_{mnu}
=
\begin{cases}
1, & \text{if UAV }u\text{ travels from waypoint }m\text{ to }n,\\
0, & \text{otherwise}.
\end{cases}
\end{equation}

The continuous variable $\tau_{mu}$ denotes the arrival time of UAV $u$ at
waypoint $m$, while $C_i$ represents the completion time of task $i$. The
variable $L_{mnu}$ represents the payload mass carried by UAV $u$ while
traversing flight link $(m,n)$. These continuous variables satisfy
\begin{equation}
\tau_{mu}\geq 0,\qquad C_i\geq 0,\qquad L_{mnu}\geq 0.
\end{equation}

\section{Time-Dependent Service Value}
\label{sec:service_value}

Each accepted survivor-assistance task $i$ is assigned an operational severity
score $s_i\in[0,1]$ by the command center, where a larger value represents
higher operational priority. PGBM treats this severity score as an input rather
than estimating it. Let $t_i$ denote the detection time of task $i$, measured
from the earthquake occurrence. The service value of the task at time $\tau$ is
defined as
\begin{equation}
v_i(\tau)
=
s_i e^{-\lambda(\tau-t_i)},
\qquad
\tau\geq t_i,\quad \lambda>0.
\end{equation}

At the detection time,
\begin{equation}
v_i(t_i)=s_i,
\end{equation}
after which the service value decreases continuously as assistance is delayed.

The parameter $\lambda$ is the common proportional decay rate; when time is
measured in minutes, its unit is $\mathrm{min}^{-1}$. The same proportional
decay rate is used for all severity levels, so severity determines the initial
task value while delay determines how that value decreases over time.
Consequently, at the same delay, a higher-severity task remains at least as
valuable as a lower-severity task. Tasks with different delays may nevertheless
have different current values, allowing the planner to account for both
severity and waiting time. The service value is an operational priority measure
and should not be interpreted as a probability of survival. The decay parameter
is calibrated experimentally in Section~\ref{sec:temporal_decay_experiment}.

\section{Optimization Problem}

The objective of the responder planner is to maximize the total
time-dependent service value obtained from the tasks served in the current
planning step:
\begin{equation}
\max
\sum_{i\in\mathcal{I}(t)}
\sum_{u\in\mathcal{U}(t)}
x_{iu}v_i(C_i).
\end{equation}

Using the service-value model defined in Section~\ref{sec:service_value},
\begin{equation}
v_i(C_i)
=
s_i
e^{-\lambda(C_i-t_i)}.
\end{equation}

A task contributes to the objective only when it is assigned to a responder
UAV. Tasks not selected in the current planning step contribute zero value and
remain pending for later consideration. Because service value decreases with
completion delay, the objective encourages the planner to serve valuable tasks
as early as possible while still respecting routing and resource constraints.
The exponential term makes the objective nonlinear, while the assignment and
routing decisions are binary. The resulting formulation is therefore a
mixed-integer nonlinear optimization problem. The constraints of the model are
listed in the following subsections as C1, C2, and so on.

\section{Joint Task Allocation and Three-Dimensional Routing}

For route accounting, the base station of UAV $u$ is written as a logical
start node $b_u^{\mathrm{s}}$ and a logical return node $b_u^{\mathrm{r}}$.
Both nodes refer to the same physical base location:

\noindent\textbf{C1: Same-base return condition.}
\begin{equation*}
\mathbf{p}_{b_u^{\mathrm{s}}}
=
\mathbf{p}_{b_u^{\mathrm{r}}},
\qquad
\forall u\in\mathcal{U}(t).
\end{equation*}

For task $i$, $m_i\in\mathcal{N}$ denotes the UAV-accessible drop-off
waypoint. This point may differ from the detected survivor coordinate. For
example, if the survivor is inside rubble or a damaged structure, the UAV may
not be able to reach the exact survivor position. In that case, $m_i$ is a
nearby feasible waypoint from which the relief parcel can be delivered or
released safely.

Each task can be assigned to at most one responder UAV:

\noindent\textbf{C2: Single-assignment condition.}
\begin{equation*}
\sum_{u\in\mathcal{U}(t)}
x_{iu}
\leq 1,
\qquad
\forall i\in\mathcal{I}(t).
\end{equation*}
If a task is not assigned in the current planning step, it remains pending and
can be considered again later.

If UAV $u$ is assigned one or more tasks, its selected route must leave its
own start-base node:

\noindent\textbf{C3: Departure-from-base condition.}
\begin{equation*}
\sum_{i\in\mathcal{I}(t)}
x_{iu}
\leq
|\mathcal{I}(t)|
\sum_{n:(b_u^{\mathrm{s}},n)\in\mathcal{E}}
y_{b_u^{\mathrm{s}},n,u},
\qquad
\forall u\in\mathcal{U}(t).
\end{equation*}

The same mission must also return to the corresponding return-base node:

\noindent\textbf{C4: Return-to-base condition.}
\begin{equation*}
\sum_{i\in\mathcal{I}(t)}
x_{iu}
\leq
|\mathcal{I}(t)|
\sum_{m:(m,b_u^{\mathrm{r}})\in\mathcal{E}}
y_{m,b_u^{\mathrm{r}},u},
\qquad
\forall u\in\mathcal{U}(t).
\end{equation*}

\section{Travel and Task Completion Time}

For each feasible flight link $(m,n)\in\mathcal{E}$, the 3D travel distance is
\begin{equation}
d_{mn}
=
\sqrt{
(p_m^x-p_n^x)^2
+
(p_m^y-p_n^y)^2
+
(p_m^z-p_n^z)^2
}.
\end{equation}

Let $\nu_u$ denote the cruise speed of UAV $u$. The corresponding travel time
on link $(m,n)$ is
\begin{equation}
t_{mnu}
=
\frac{d_{mn}}{\nu_u}.
\end{equation}

The total selected travel distance and travel time for UAV $u$ are
\begin{equation}
D_u
=
\sum_{(m,n)\in\mathcal{E}}
d_{mn}y_{mnu},
\end{equation}
and
\begin{equation}
T_u
=
\sum_{(m,n)\in\mathcal{E}}
t_{mnu}y_{mnu}.
\end{equation}
These sums include the departure from base, travel between selected drop-off
waypoints, and the return-to-base link.

Along a selected flight link, arrival time must increase by at least the
corresponding travel time:

\noindent\textbf{C5: Selected-link timing condition.}
\begin{equation*}
y_{mnu}=1
\Rightarrow
\tau_{nu}
\geq
\tau_{mu}+t_{mnu}.
\end{equation*}

If task $i$ is assigned to UAV $u$, its completion time is the arrival time at
the drop-off waypoint plus the service duration:

\noindent\textbf{C6: Task-completion time condition.}
\begin{equation*}
C_i
=
\tau_{m_i u}+\delta_i,
\qquad
\text{if }x_{iu}=1.
\end{equation*}

\section{Relief Capacity and Payload Evolution}

Let $\mathcal{R}$ be the set of relief-item types. At time $t$, task $i$
requires $q_{ir}(t)$ remaining items of type $r$, and each item of type $r$
has weight $w_r$. UAV $u$ can carry at most $K_u$ parcels:

\noindent\textbf{C7: Parcel-count capacity condition.}
\begin{equation*}
\sum_{i\in\mathcal{I}(t)}
\sum_{r\in\mathcal{R}}
q_{ir}(t)x_{iu}
\leq
K_u.
\end{equation*}

The same selected load must also satisfy the payload-weight capacity:

\noindent\textbf{C8: Payload-weight capacity condition.}
\begin{equation*}
\sum_{i\in\mathcal{I}(t)}
\sum_{r\in\mathcal{R}}
w_r q_{ir}(t)x_{iu}
\leq
W_u.
\end{equation*}

The relief mass required by task $i$ is
\begin{equation}
M_i(t)
=
\sum_{r\in\mathcal{R}}
w_r q_{ir}(t),
\end{equation}
and the initial payload assigned to UAV $u$ is
\begin{equation}
M_u^0
=
\sum_{i\in\mathcal{I}(t)}
M_i(t)x_{iu}.
\end{equation}

The payload carried on a flight link can be positive only if that link is
selected:

\noindent\textbf{C9: Selected-link payload condition.}
\begin{equation*}
L_{mnu}
\leq
W_u y_{mnu}.
\end{equation*}

At departure, the UAV carries the complete relief mass assigned to its
mission:

\noindent\textbf{C10: Initial payload condition.}
\begin{equation*}
\sum_{n:(b_u^{\mathrm{s}},n)\in\mathcal{E}}
L_{b_u^{\mathrm{s}},n,u}
=
M_u^0.
\end{equation*}

After task $i$ is served, the carried payload decreases by the mass delivered
at that drop-off waypoint:

\noindent\textbf{C11: Payload-reduction condition.}
\begin{equation*}
\sum_{n:(m_i,n)\in\mathcal{E}}
L_{m_i n u}
=
\sum_{m:(m,m_i)\in\mathcal{E}}
L_{m m_i u}
-
M_i(t)x_{iu}.
\end{equation*}

\section{Energy and Battery Feasibility}

Battery feasibility is checked against the available energy $E_u(t)$ while
keeping reserve energy $E_u^{\mathrm{res}}$ unused for planned service. For
each feasible flight link $(m,n)$, positive ascent is
\begin{equation}
a_{mn}
=
\max(0,p_n^z-p_m^z).
\end{equation}

The required mission energy is modeled as
\begin{equation}
\begin{aligned}
E_u^{\mathrm{req}}
=&
\sum_{(m,n)\in\mathcal{E}}
\alpha_u d_{mn}y_{mnu}
+
\sum_{(m,n)\in\mathcal{E}}
\beta_u d_{mn}L_{mnu}
\\
&+
\sum_{(m,n)\in\mathcal{E}}
\gamma_u a_{mn}y_{mnu}
+
\sum_{i\in\mathcal{I}(t)}
P_u^{\mathrm{srv}}\delta_i x_{iu}.
\end{aligned}
\end{equation}

The four terms represent baseline travel energy, payload-related travel
energy, positive-ascent energy, and service or hovering energy. Since the
selected route includes the return-to-base link, the energy required for safe
return is included in the same expression.

The mission is feasible only if the required energy is no greater than the
usable battery energy:

\noindent\textbf{C12: Energy-feasibility condition.}
\begin{equation*}
E_u^{\mathrm{req}}
\leq
E_u(t)-E_u^{\mathrm{res}}.
\end{equation*}
The reserve term leaves a safety margin for operational uncertainty and safe
recovery.

\section{Task-Replacement Recourse}

New survivor-assistance tasks may arrive while all responder UAVs are already
deployed. In that case, the planner does not restart the full mission plan.
Instead, it checks whether the newly arrived task should replace one
uncompleted task in an active mission.

Let $n$ denote the new task and let $t_r$ denote the recourse decision time.
For UAV $u$, define
\[
\mathcal{T}_u^{\mathrm{rem}}(t_r)
\]
as the set of tasks that remain uncompleted at time $t_r$. Each candidate
recourse action replaces one task
$i\in\mathcal{T}_u^{\mathrm{rem}}(t_r)$ with the newly arrived task $n$. The
completed part of the mission remains fixed.

Before replacement, the value of the remaining mission is
\begin{equation}
J_{ui}^{\mathrm{old}}
=
\sum_{j\in\mathcal{T}_u^{\mathrm{rem}}(t_r)}
v_j(C_j^{\mathrm{old}}).
\end{equation}

After task $i$ is replaced by task $n$, the remaining route is recalculated.
The new remaining value is
\begin{equation}
J_{ui}^{\mathrm{new}}
=
\sum_{j\in(\mathcal{T}_u^{\mathrm{rem}}(t_r)\setminus\{i\})\cup\{n\}}
v_j(C_j^{\mathrm{new}}).
\end{equation}

The gain from this replacement is
\begin{equation}
\Delta J_{ui}
=
J_{ui}^{\mathrm{new}}
-
J_{ui}^{\mathrm{old}}.
\end{equation}

A replacement is feasible only if UAV $u$ already carries the relief items
required by the new task:
\begin{equation}
q_{nr}(t_r)
\leq
I_{ur}(t_r),
\qquad
\forall r\in\mathcal{R}.
\end{equation}
The modified remaining mission must also satisfy battery feasibility:
\begin{equation}
E_{ui}^{\mathrm{new}}
\leq
E_u(t_r)-E_u^{\mathrm{res}},
\end{equation}
and its recalculated route must use feasible flight links in $\mathcal{E}$ and
return safely to the UAV base.

Let $\mathcal{S}(t_r)$ be the set of feasible one-for-one replacements at
time $t_r$. The selected replacement is
\begin{equation}
(u^*,i^*)
=
\arg\max_{(u,i)\in\mathcal{S}(t_r)}
\left\{
\Delta J_{ui}
\right\}.
\end{equation}

The replacement is accepted only if it produces a positive gain:
\begin{equation}
\Delta J_{u^*i^*}>0.
\end{equation}
If accepted, task $n$ is inserted into the remaining mission of UAV $u^*$, and
the displaced task $i^*$ is returned to the waiting queue:
\begin{equation}
\mathcal{Q}(t_r^+)
=
\mathcal{Q}(t_r)\cup\{i^*\}.
\end{equation}

If no feasible replacement produces a positive gain, the active missions
remain unchanged and the newly arrived task is added to the waiting queue:
\begin{equation}
\mathcal{Q}(t_r^+)
=
\mathcal{Q}(t_r)\cup\{n\}.
\end{equation}

This recourse rule is intentionally limited. It updates one active mission
only when doing so improves the remaining service value, and it does not
invalidate route segments or deliveries that have already been completed.
